\subsection{Lectura de los resultados}

\begin{observacion}[El mecanismo distribuido entrega, y con qué tasa]
\label{obs:ensayo-cierre-tasa}
Con la configuración de referencia —Smith, aptitud vector, arrendamiento,
seguridad anidada, reclutamiento atómico, creencias exactas— las dos cargas se
entregan en $21$ de $30$ mundos, con intervalo de Wilson del $52$ al $83\%$, en
una mediana de $153$~s desde el instante cero hasta la segunda entrega, con un
reclutamiento que converge en unas seis revisiones aceptadas. Los nueve fallos
son de dos tipos: formaciones que no cierran antes de los sesenta segundos de
guarda y se reabren (churn), y segundas cargas que llegan al horizonte en
transporte. Ninguno es un certificado falso: en las trescientas misiones de la
rejilla de protocolos no hubo una sola coalición certificada con margen negativo.
Es la primera vez en este documento que el mecanismo completo —revisión,
certificado, mercado, arrendamiento, barreras— corre solo sobre la planta; el
$70\%$ no es la cifra de un sistema industrial, es la de una construcción
teórica ejecutada sin un solo nodo con el estado completo.
\end{observacion}

\begin{observacion}[Los protocolos con estacionariedad de Nash son equivalentes; el replicador no recluta]
\label{obs:ensayo-cierre-protocolos}
Mejor respuesta, Smith, BNN y logit entregan entre $19$ y $22$ de $30$, con
McNemar pareado frente a Smith de $p\ge0{,}8$ en los tres casos: sobre esta
planta el protocolo de revisión no cambia el resultado, como predice
\cref{obs:smith} para los protocolos que comparten estacionariedad de Nash. El
replicador entrega $0$ de $30$ con cero revisiones aceptadas en trescientas
misiones-segundo: es la invariancia de la frontera del símplex medida —una
estrategia que nadie usa no puede ser imitada— y frente a Smith da $0/21$
discordancias, $p<10^{-6}$. \Cref{obs:smith} lo enunció como razón para
descartar el replicador en reclutamiento; el banco lo convierte en un hecho
con intervalo.
\end{observacion}

\begin{observacion}[La aptitud vector no fue vinculante en \textsf{cargo}]
\label{obs:ensayo-cierre-aptitud}
Aptitud vector y escalar dan trayectorias idénticas en las treinta semillas:
mismo éxito, mismos tiempos, mismos mensajes. No es coincidencia sino
geometría: con tres pads de disco bajo una plataforma, toda tripleta con
capacidad suficiente tiene también margen de \emph{wrench} en todas las
direcciones, luego el certificado de \cref{th:span} y la suma de capacidades
certifican las mismas coaliciones. La diferencia entre ambas aptitudes es la
de \cref{th:span} frente al rango: aparece cuando las columnas son
unilaterales —empuje— y ese es exactamente el modo que no entró en el banco.
El resultado es, por tanto, un control negativo honesto: dice dónde la
aptitud vector no añade nada, que es donde la teoría dice que no debe añadir.
\end{observacion}

\begin{observacion}[Arrendamiento y precio de congestión: misma tasa, distinto coste]
\label{obs:ensayo-cierre-pasillo}
El arrendamiento exclusivo con orden total entrega $15$ de $20$ y el precio de
congestión sin exclusividad $14$ de $20$, con una sola discordancia pareada:
en tasa de éxito no se distinguen. Se distinguen en lo que cuestan: el precio
consume un $22\%$ más de mensajes y un $20\%$ más de energía, con una mediana
de misión siete segundos más larga, porque las dos coaliciones renegocian
velocidad en cada paso de control mientras se aproximan al hueco en lugar de
esperar en el punto de retención. En ninguna de las cuarenta misiones dos
centros de carga ocuparon el vano a la vez: en esta geometría —un hueco de
$3{,}3$~m por el que no caben dos plataformas de $1{,}1$ y $1{,}3$~m con sus
AMR— la exclusión la impone la física antes que el contrato, y el precio la
reproduce sin garantía. Es la tercera lectura de \cref{obs:ensayo-cuestiona} y
de \cref{obs:mixed2d-exclusividad}: la garantía de \cref{th:ciclo-espera} no
sale gratis, pero aquí su precio en tiempo es nulo y en comunicación es
negativo.
\end{observacion}

\begin{observacion}[La capa de seguridad, tal como está implementada, cuesta éxito y no elimina el contacto]
\label{obs:ensayo-cierre-seguridad}
Es el resultado incómodo del banco y se reporta sin atenuarlo. Sin barreras
—solo la repulsión de contacto de la planta— el mecanismo entrega $19$ de
$20$ en una mediana de $110$~s; con la barrera de coalición sola, $18$ de $20$
en $144$~s; con la anidada (coalición más AMR), $15$ de $20$ en $153$~s. Las
diferencias pareadas frente a la anidada son $5/1$ y $5/2$ discordancias,
$p=0{,}22$ y $p=0{,}45$: no significativas con veinte mundos, pero de un solo
signo. Y la separación mínima entre AMR no acoplados es negativa en las tres
variantes —$-0{,}13$, $-0{,}16$ y $-0{,}22$~m de mediana—, es decir, hay
contacto entre parachoques en casi todas las misiones; la barrera anidada lo
reduce, no lo elimina. La lectura correcta es sobre \cref{th:entrega}: la
hipótesis (A5) exige barreras \emph{factibles en todo estado alcanzado}, y lo
implementado es un límite heurístico de velocidad por distancia de frenado
sobre la navegación de robots que además deben pasar bajo la plataforma. Esa
heurística frena a los que se acercan a formar —de ahí las formaciones que no
cierran— y no impide que dos robots que se rodean acaben tocándose. El
teorema no falla; falla la hipótesis, y el banco dice cuál: hace falta el
programa cuadrático con barrera de orden correcto de
\cref{th:grado-relativo-rueda} y garantía de factibilidad, no un límite de
velocidad. Hasta entonces, el sistema que entrega es el que no frena, y eso
es exactamente lo que \cref{obs:ensayo-cuestiona} ya había medido en otra
planta.
\end{observacion}

\begin{observacion}[Reclutamiento atómico y GNE primal--dual: mismo resultado, más mensajes]
\label{obs:ensayo-cierre-gne}
Sustituir la revisión atómica por dinámicas de Smith sobre cuotas continuas
con multiplicadores de capacidad y atomizar a los doce segundos entrega $16$
de $20$ frente a $15$ de $20$, con $4/3$ discordancias pareadas ($p=1$): la
relajación continua y su redondeo llegan a las mismas coaliciones que la
revisión atómica, con la mitad de revisiones aceptadas —tres frente a seis,
porque la atomización ya nace cerca de un equilibrio— y un $35\%$ más de
mensajes, porque las cuotas se difunden en cada paso de control mientras se
relajan. Es \cref{th:vi} y \cref{obs:vi} medidos: el equilibrio variacional
de la relajación es alcanzable de forma distribuida, y su brecha de
integralidad frente a la revisión atómica es, en estas instancias, nula en
resultado y positiva en comunicación. Ninguna de las dos vías produjo un
certificado falso.
\end{observacion}

\begin{observacion}[Hiperjuego con margen holgado: las creencias erróneas no cambian nada]
\label{obs:ensayo-cierre-hiperjuego}
Con error uniforme de $\pm35\%$ en las capacidades ajenas y en las masas, con
o sin recertificación, las veinte semillas producen trayectorias idénticas a
las de creencias exactas, y no hay un solo certificado falso. La razón es la
misma que en \cref{obs:ensayo-cierre-aptitud}: en \textsf{cargo} la demanda
de diseño —unos $30$~N para $60$~kg a $0{,}35$~m/s$^2$— está a un factor
tres o más de la capacidad de cualquier tripleta, y un error del $35\%$ no
cruza esa distancia. El hiperjuego solo muerde cuando el certificado es
vinculante; por eso la rejilla siguiente reduce las capacidades a un $40\%$ y
eleva la aceleración de diseño a $1$~m/s$^2$, que es donde la recertificación
de \cref{obs:certificado-postevento} tiene algo que rechazar.
\end{observacion}

\begin{observacion}[Qué queda demostrado y ensayado de la tesis central, y qué no]
\label{obs:ensayo-cierre-balance}
La tesis central de este documento es que reclutamiento, modo, contacto y ruta
pueden decidirse como un solo juego distribuido de $n$ saltos sobre robots
diferenciales heterogéneos con contacto real, y que ese juego entrega la
carga. \Cref{th:potencial-valor} y \cref{th:entrega} la formulan y la
demuestran bajo hipótesis nombradas; el banco la ensaya. Lo que el banco
sostiene: el mecanismo completo, sin ningún nodo con el estado global, entrega
las dos cargas en el $63$--$73\%$ de los mundos según el protocolo, con cero
certificados falsos en $430$ misiones, recluta en seis revisiones, cruza el
hueco por arrendamiento sin doble ocupación, y su capa continua no depende del
protocolo de revisión ni de la relajación elegida. Lo que el banco niega: que
el replicador recluta (nunca lo hace), que la aptitud vector importe en
\textsf{cargo} con pads de disco (no importa), que un error de creencia del
$35\%$ importe con margen holgado (no importa), y que la barrera heurística
implementada cumpla (A5) (no la cumple: hay contacto entre parachoques en
casi todas las misiones y la barrera cuesta éxito). Lo que el banco no toca:
el empuje unilateral, la planta tridimensional, el hardware, y la optimalidad.
Tres de cada diez misiones fallan por formaciones que no cierran o por
segundas cargas que agotan el horizonte; ese $30\%$ es el número que una
versión futura tiene que bajar, y \cref{obs:ensayo-cierre-seguridad} dice
por dónde empezar.
\end{observacion}

\begin{observacion}[Hiperjuego con margen vinculante: la recertificación hace su trabajo]
\label{obs:ensayo-cierre-hiperjuego2}
Con capacidades al $40\%$ y aceleración de diseño de $1$~m/s$^2$, el
certificado pasa a ser vinculante y el error de creencias del $\pm35\%$
produce lo que la teoría anuncia. Sin recertificación, tres coaliciones se
certifican con creencias que no lo estaban con las capacidades reales, entran
en formación con ese margen negativo y la tasa baja a $14$ de $20$; con
recertificación al formar, el mecanismo detecta y rechaza el certificado
falso veinticinco veces en veinte mundos —expulsa al miembro de menor
capacidad y reabre el reclutamiento— y la tasa vuelve a $15$ de $20$, la
misma que con creencias exactas, sin un solo certificado falso que llegue al
transporte. La diferencia pareada entre recertificar y no hacerlo es de una
sola misión ($p=1$ con veinte mundos), así que lo que el banco establece no
es una ventaja de tasa sino un mecanismo: bajo creencias erróneas y margen
vinculante, la recertificación con capacidades observadas es lo que impide
que el juego apruebe una continuación con capacidades que no existen, que es
\cref{obs:certificado-postevento} ejecutado dentro del mecanismo distribuido
y no en un frenado aislado. El coste es churn: veinticinco reaperturas, con
el mismo número de mensajes y la misma energía.
\end{observacion}

\begin{observacion}[El juego de trayectorias frente al seguidor: más éxito y menos contacto, a cambio de tiempo y mensajes]
\label{obs:ensayo-cierre-planificador}
Sobre las mismas treinta semillas, el juego continuo de trayectorias entrega
las dos cargas en $24$ de $30$ frente a $20$ de $30$ del seguidor de puntos,
con $7/4$ discordancias pareadas ($p=0{,}55$: no significativo a este tamaño)
y una separación mínima mediana entre robots de $-0{,}12$ frente a
$-0{,}16$~m; cuesta un $8\%$ más de tiempo, un $12\%$ más de energía y un
$13\%$ más de mensajes, porque cada jugador publica su continuación cinco veces
por segundo. En la semilla de referencia el cambio de rumbo por metro recorrido
de las cargas baja de $3{,}0$ y $2{,}3$ a $2{,}0$ y $1{,}5$~rad/m —un tercio
menos de ondulación— mientras la longitud recorrida sube para la carga que
espera, por el desvío al carril. Las rejillas principales de esta sección se
ejecutaron con el seguidor de puntos, que es la capa con la que se obtuvieron
sus tablas; la comparación directa entre ambas capas es esta fila, y la capa
de juego es la que se conserva como predeterminada. Lo que el resultado no
dice es que las rutas sean óptimas: son un equilibrio local del potencial de
\eqref{eq:potencial-trayectorias} sobre una ruta discreta impuesta, con
desvíos que la geometría de los carriles produce, y así se ve en
\cref{fig:sim-plan}.
\end{observacion}
